% Source-derived chapter 03: The Geometric Syzygy Conjecture
% Original source: universal-secant-bundles-syzygies-canonical-curves
\chapter{The Geometric Syzygy Conjecture}
\label{universal-secant-bundles-syzygies-canonical-curves-chapter-03}
\source{universal-secant-bundles-syzygies-canonical-curves}

\begin{remark}[Source section]
\label{universal-secant-bundles-syzygies-canonical-curves-chapter-03-source}
\source{universal-secant-bundles-syzygies-canonical-curves}
\source{universal-secant-bundles-syzygies-canonical-curves}
This chapter reproduces the corresponding section of the retrieved
LaTeX source, including its definitions, statements, examples, and
proofs. It has not yet been translated into Lean.
\notready
\end{remark}

% The source section title was: The Geometric Syzygy Conjecture
\label{sec2}
\source{universal-secant-bundles-syzygies-canonical-curves}
In this section, we use the techniques used in our proof of Voisin's Theorem to resolve the Geometric Syzygy Conjecture for extremal syzygies of generic curves of even genus. We stick with the notation from Section \ref{sec1}.  As before, consider the exact sequence $0 \to \mathcal{S} \to \pi^* \mathcal{M} \to \Gamma \to 0.$
\begin{thm}
\source{universal-secant-bundles-syzygies-canonical-curves}
\notready \label{real-thm}
\source{universal-secant-bundles-syzygies-canonical-curves}
The natural morphism $\psi: \mathrm{H}^1(B, \wedge^k \pi^* \mathcal{M} \otimes {p'}^*L) \to \mathrm{H}^1(B, \wedge^k \Gamma \otimes {p'}^*L)$ is surjective.
\end{thm}
\begin{proof}
From the exact sequence
\begin{align*}
\ldots \to \mathcal{S} \otimes \wedge^{k-1} \pi^* \mathcal{M}   \otimes {p'}^*L \to \wedge^k  \mathcal{M}   \otimes {p'}^*L \to \wedge^k \Gamma \otimes {p'}^*L \to 0,
\end{align*}
it suffices to show $\mathrm{H}^{1+i}(\mathrm{Sym}^i \mathcal{S} \otimes \wedge^{k-i} \pi^* \mathcal{M} \otimes {p'}^*L)=0$
for $1 \leq i \leq k$. From the exact sequence
$$0 \to \mathrm{Sym}^{i} \mathcal{S} \to \mathrm{Sym}^i( \pi^*\mathcal{G}) \to \mathrm{Sym}^{i-1}(\pi^*\mathcal{G})\otimes {p'}^*L \otimes I_D \to 0$$
it suffices to have
%\begin{align*}
%&\mathrm{H}^{1+i}\left(B,\mathrm{Sym}^i\left({q'}^*q'_*({p'}^*L \otimes I_D)\right) \otimes \wedge^{k-i} \pi^* \mathcal{M} \otimes {p'}^*L\right)=0,\\
%&\mathrm{H}^i\left(B,\mathrm{Sym}^{i-1}\left({q'}^*q'_*({p'}^*L \otimes I_D)\right)\otimes \wedge^{k-i} \pi^* \mathcal{M} \otimes {p'}^*L^{\otimes 2} \otimes I_D\right)=0
%\end{align*}q
%This is equivalent to 
\begin{align*}
\mathrm{H}^{1+i}\left(X \times \PP,\mathrm{Sym}^i (\mathcal{G}) \otimes \wedge^{k-i} \mathcal{M} \otimes {p}^*L\right)=\mathrm{H}^i\left(X \times \PP,\mathrm{Sym}^{i-1}(\mathcal{G})\otimes \wedge^{k-i}  \mathcal{M} \otimes {p}^*L^{\otimes 2} \otimes I_{\mathcal{Z}}\right)=0
\end{align*}
By taking $\mathrm{Sym}^i$ of the exact sequence
$$0 \to q^*\mathcal{O}_{\PP}(-2) \to \mathrm{H}^0(E) \otimes q^*\mathcal{O}_{\PP}(-1) \to \mathcal{G} \to 0$$
it suffices to show the four vanishings
{\small{\begin{align*}
&\mathrm{H}^{2+i}\left( \wedge^{k-i} M_L (L) \boxtimes   \mathcal{O}_{\PP}(-i-1)  \right)=0, & &\mathrm{H}^{1+i}\left(\wedge^{k-i} M_L (L) \boxtimes \mathcal{O}_{\PP}(-i)  \right)=0,\\
&\mathrm{H}^{1+i}\left((\wedge^{k-i} M_L (2L) \boxtimes \mathcal{O}_{\PP}(-i) )\otimes I_{\mathcal{Z}} \right)=0, & &\mathrm{H}^{i}\left((\wedge^{k-i} M_L (2L) \boxtimes \mathcal{O}_{\PP}(-i+1) )\otimes I_{\mathcal{Z}} \right)=0.
\end{align*}}}
The first two claims follow from the K\"unneth formula, for $1 \leq i \leq k$. For the last two claims, we use the short exact sequence
$$0 \to L^{-1} \boxtimes \mathcal{O}_{\PP}(-2) \to E(L^{-1}) \boxtimes \mathcal{O}_{\PP}(-1) \to I_{\mathcal{Z}} \to 0,$$
so it suffices to have the four vanishings
{\small{\begin{align*}
&\mathrm{H}^{2+i}\left(\wedge^{k-i} M_L (L) \boxtimes \mathcal{O}_{\PP}(-i-2)  \right)=0,
& &\mathrm{H}^{1+i}\left( \wedge^{k-i} M_L \otimes E(L) \boxtimes \mathcal{O}_{\PP}(-i-1)  \right)=0,\\
&\mathrm{H}^{1+i}\left( \wedge^{k-i} M_L (L) \boxtimes \mathcal{O}_{\PP}(-i-1)  \right)=0,
& &\mathrm{H}^{i}\left(\wedge^{k-i} M_L \otimes E(L) \boxtimes \mathcal{O}_{\PP}(-i)  \right)=0.
\end{align*}}}
This follows from the K\"unneth formula, using $\mathrm{H}^1(X,L)=0$ if $i=k$ in the first vanishing.
\end{proof}

The Geometric Syzygy Conjecture in even genus now follows readily from Theorem \ref{real-thm}. 
\begin{lem}
\source{universal-secant-bundles-syzygies-canonical-curves}
\notready
With notation as in Section \ref{sec1}, we have $\wedge^k \Gamma \simeq I_D \otimes {q'}^* \mathcal{O}_{\PP}(k)$.
\end{lem}
\begin{proof}
By taking determinants of the exact sequence $0 \to \Gamma \to {q'}^* \mathcal{W} \to {p'}^* L_{|_D} \to 0$, it suffices to show $\det \mathcal{W} \simeq \mathcal{O}_{\PP}(k)$. From $0 \to q_*(p^*L \otimes I_{\mathcal{Z}}) \to q_*p^*L \to \mathcal{W} \to 0$, we see $\det \mathcal{W}=\det (q_*(p^*L \otimes I_{\mathcal{Z}}))^{-1}$. We now deduce the claim from the exact sequence
$$0 \to \mathcal{O}_{\PP}(-2) \to \mathrm{H}^0(E) \otimes \mathcal{O}_{\PP}(-1) \to q_*(p^*L \otimes I_{\mathcal{Z}}) \to 0.$$\end{proof}
\begin{cor}
\source{universal-secant-bundles-syzygies-canonical-curves}
\notready \label{main-cor-nat-iso}
\source{universal-secant-bundles-syzygies-canonical-curves}
The morphism $\psi: \mathrm{H}^1(\wedge^k \pi^* \mathcal{M} \otimes {p'}^*L) \to \mathrm{H}^1(\wedge^k \Gamma \otimes {p'}^*L)$ is an isomorphism and induces a natural isomorphism $\mathrm{K}_{k-1,1}(X,L) \simeq \mathrm{Sym}^{k-2}\mathrm{H}^0(X,E)  $.
\end{cor}
\begin{proof}
We have $\mathrm{H}^1(B, I_D \otimes {q'}^* \mathcal{O}_{\PP}(k))\simeq \mathrm{H}^1(X \times \PP, I_{\mathcal{Z}} \otimes {q}^* \mathcal{O}_{\PP}(k))$. From
$$0 \to \mathcal{O}_X \boxtimes \mathcal{O}_{\PP}(-2) \to E \boxtimes \mathcal{O}_{\PP}(-1) \to p^*L \otimes  I_{\mathcal{Z}} \to 0,$$
we obtain an isomorphism $\mathrm{H}^1(B, \wedge^k \Gamma \otimes {p'}^*L) \simeq \mathrm{H}^0(\mathcal{O}_{\PP}(k-2)) \simeq {\mathrm{Sym}^{k-2}\mathrm{H}^0(X,E)}^{\vee}$. Theorem \ref{real-thm} gives a surjective map $\mathrm{K}_{k-1,2}(X,L) \to {\mathrm{Sym}^{k-2}\mathrm{H}^0(X,E)}^{\vee}$. Dualizing, we have an injective map $\mathrm{Sym}^{k-2}\mathrm{H}^0(X,E) \to \mathrm{K}_{k-1,1}(X,L).$ Voisin's Theorem in even genus, as proven in Section \ref{sec1}, implies $\dim \mathrm{Sym}^{k-2}\mathrm{H}^0(X,E) = \dim \mathrm{K}_{k-1,1}(X,L)$,  \cite[\S 4.1]{farkas-progress}, completing the proof.\end{proof}

We record an observation for future use. From the sequence $0 \to \mathcal{S} \to \pi^* \mathcal{M} \to \Gamma \to 0$, the isomorphism $\psi: \mathrm{H}^1( \wedge^k \pi^* \mathcal{M} \otimes {p'}^*L) \xrightarrow{\sim} \mathrm{H}^1(\wedge^k \Gamma \otimes {p'}^*L)$ is Serre dual to
$$\psi^{\vee} \; : \; \mathrm{H}^{k+2}( \wedge^k \mathcal{S} \otimes \omega_B) \to \mathrm{H}^{k+2}( \wedge^k \pi^* \mathcal{M}\otimes \omega_B).$$
%It is possible to unravel the identifications in the isomorphism $\mathrm{H}^1( \wedge^k \pi^* \mathcal{M} \otimes {p'}^*L) \xrightarrow{\sim} \mathrm{H}^1(\wedge^k \Gamma \otimes {p'}^*L)$.
\begin{proof}[Proof of Theorem \ref{geo-thm-intro}]
We have the exact sequence $0 \to \Gamma' \to q^* \mathcal{W} \to p^*L_{|_{\mathcal{Z}}} \to 0,$ where we have set $\Gamma':=\pi_* \Gamma$. The image of the natural map $\wedge^k \Gamma' \to \wedge^k q^* \mathcal{W}  \simeq \mathcal{O}_{\PP}(k)$ is $I_{\mathcal{Z}} \otimes q^*\mathcal{O}_{\PP}(k)$ from \cite[Cor.\ 3.7]{ein-lazarsfeld-asymptotic}. The induced surjection $\wedge^k \Gamma' \to I_{\mathcal{Z}} \otimes q^*\mathcal{O}_{\PP}(k)$ can be identified with the natural map
$$\wedge^k \pi_* \Gamma \to \pi_*\wedge^k \Gamma \simeq \pi_*(I_D \otimes {q'}^*\mathcal{O}_{\PP}(k)) \hookrightarrow \pi_* \wedge^k {q'}^* \mathcal{W} \simeq \wedge^k q^* \mathcal{W} . $$ The map $\psi: \mathrm{H}^1(B, \wedge^k \pi^* \mathcal{M} \otimes {p'}^*L) \to \mathrm{H}^1(B, \wedge^k \Gamma \otimes {p'}^*L)$ can be identified with the composition
$$\mathrm{H}^1(X\times \PP, \wedge^k\mathcal{M} \otimes p^*L) \to \mathrm{H}^1(X\times \PP, \wedge^k \Gamma' \otimes p^*L) \to \mathrm{H}^1( X\times \PP, (L \boxtimes \mathcal{O}_{\PP}(k))\otimes I_{\mathcal{Z}}),$$
where the first map is induced from the natural surjection $\mathcal{M} \to \Gamma'$. Since $q_*\wedge^k\mathcal{M} \otimes p^*L \simeq \mathrm{H}^0(X, \wedge^k M_L(L)) \otimes \mathcal{O}_{\PP}$, we have $\mathrm{H}^i(q_*\wedge^k\mathcal{M} \otimes p^*L)=0$ for $i>0$. Thus the Leray spectral sequence gives an isomorphism
$$\mathrm{H}^1(X \times \PP, \wedge^k\mathcal{M} \otimes p^*L) \simeq \mathrm{H}^0(\PP, R^1q_*\wedge^k\mathcal{M} \otimes p^*L).$$
Next, the exact sequence $0 \to \mathcal{O}_{\PP}(k-2) \to \mathrm{H}^0(E) \otimes \mathcal{O}_{\PP}(k-1) \to q_*((L \boxtimes \mathcal{O}_{\PP}(k))\otimes I_{\mathcal{Z}}) \to 0$ gives $\mathrm{H}^i(q_*((L \boxtimes \mathcal{O}_{\PP}(k))\otimes I_{\mathcal{Z}}))=0$ for $i >0$ so 
$$\mathrm{H}^1( X \times \PP,(L \boxtimes \mathcal{O}_{\PP}(k))\otimes I_{\mathcal{Z}})\simeq \mathrm{H}^0(\PP, R^1 q_*((L \boxtimes \mathcal{O}_{\PP}(k))\otimes I_{\mathcal{Z}})).$$
The map $\psi$ is thus naturally identified with the global section map applied to the morphism
$$\mathrm{K}_{k-1,1}(X,L)^{\vee} \otimes \mathcal{O}_{\PP} \simeq R^1q_*\wedge^k\mathcal{M} \otimes p^*L \to R^1 q_*((L \boxtimes \mathcal{O}_{\PP}(k))\otimes I_{\mathcal{Z}}),$$
of vector bundles. Note that $R^1 q_*((L \boxtimes \mathcal{O}_{\PP}(k))\otimes I_{\mathcal{Z}})$ is a line bundle from the proof of Lemma \ref{very-first-lem}. From
$0 \to \mathcal{O}_X \boxtimes \mathcal{O}_{\PP}(-2) \to E \boxtimes \mathcal{O}_{\PP}(-1) \to p^*L \otimes  I_{\mathcal{Z}} \to 0,$
we have $$R^1 q_*((L \boxtimes \mathcal{O}_{\PP}(k))\otimes I_{\mathcal{Z}}) \simeq R^2q_*\mathcal{O}_{X \times \PP} \otimes \mathcal{O}_{\PP}(k-2)\simeq \mathcal{O}_{\PP}(k-2),$$
using that $R^2q_*\mathcal{O}_{X \times \PP} \simeq \mathcal{O}_{\PP}$ by relative duality. Since the induced morphism
\begin{align} \label{bundle-ind-vero}
\source{universal-secant-bundles-syzygies-canonical-curves}
\mathrm{K}_{k-1,1}(X,L)^{\vee} \otimes \mathcal{O}_{\PP} \to R^1 q_*((L \boxtimes \mathcal{O}_{\PP}(k))\otimes I_{\mathcal{Z}}) \simeq \mathcal{O}_{\PP}(k-2)
\end{align}
is surjective on global sections and $\mathcal{O}_{\PP}(k-2)$ is globally generated, it must be surjective as a morphism of sheaves. Morphism (\ref{bundle-ind-vero}) therefore induces a Veronese morphism $$\widetilde{\psi} : \PP(\mathrm{H}^0(E)) \to \PP(\mathrm{K}_{k-1,1}(X,L))$$ of degree $k-2$.\smallskip

To complete the proof, it only remains to show that, for any nonzero $t \in \mathrm{H}^0(E)$, we have
$$\widetilde{\psi}([t])=[\mathrm{K}_{k-1,1}\left(X,L,\mathrm{H}^0(L \otimes I_{Z(t)})\right)].$$
The fiber of the morphism (\ref{bundle-ind-vero}) over $[t]$ is identified with the natural composition
$$\mathrm{H}^1(X,\wedge^k M_L(L)) \to \mathrm{H}^1(X,\wedge^k \Gamma'_t(L)) \to \mathrm{H}^1(X,L \otimes I_{Z(t)}),$$
using the notation from Section \ref{sec1}. As above, this may be identified with the natural surjection
$$\mathrm{H}^1(B_t,\wedge^k \pi^*_tM_L(L)) \twoheadrightarrow \mathrm{H}^1(B_t,\wedge^k \Gamma_t (\pi^*_t L )).$$
From the exact sequence $0 \to \mathcal{S}_t \to \pi_t^* M_L \to \Gamma_t \to 0$ and the isomorphisms $\wedge^k M_L(L) \simeq \wedge^k M_L^{\vee}$, $\wedge^k \Gamma_t (\pi^*_t L) \simeq \wedge^k \mathcal{S}_t^{\vee}$, we may identify the fiber of (\ref{bundle-ind-vero}) with the map $\mathrm{H}^1(X,\wedge^k M_L^{\vee}) \twoheadrightarrow \mathrm{H}^1(B_t, \wedge^k \mathcal{S}_t^{\vee})$, which is Serre dual to
$$ \mathrm{H}^1(B_t, \wedge^k \mathcal{S}_t(D_t)) \hookrightarrow \mathrm{H}^1(B_t,\wedge^k \pi_t^* M_L(D_t)). $$ Using the Kernel Bundle description of syzygies, this may be interpreted as an inclusion
$$\mathrm{K}_{k-1,1}(B_t,D_t, \pi_t^*L-D_t)\hookrightarrow \mathrm{K}_{k-1,1}(B_t,D_t,\pi_t^*L).$$
Since $\mathrm{H}^0(B_t, \pi_t^*L^{\otimes n}(D_t)) \simeq \mathrm{H}^0(B_t, \pi_t^*L^{\otimes n})\simeq \mathrm{H}^0(X, L^{\otimes n})$ for any $n$,
this inclusion is naturally identified with $\mathrm{K}_{k-1,1}\left(X,L,\mathrm{H}^0(L \otimes I_{Z(t)})\right) \hookrightarrow \mathrm{K}_{k-1,1}(X,L),$
as required.
\end{proof}

It is well-known that Theorem \ref{geo-thm-intro}, combined with an observation of Voisin, implies Corollary \ref{geo-cor}, see the unpublished \cite[\S 11]{bothmer-preprint}. For the convenience of the reader we provide the proof.
\begin{proof}[Proof of Corollary \ref{geo-cor}]
Let $X$ be as in Theorem \ref{geo-thm-intro} and let $C \in |L|$ be general. We have a surjective, finite morphism
$$d: \mathrm{Gr}_2(\mathrm{H}^0(E)) \to |L|$$
which takes a two dimensional subspace $W \seq \mathrm{H}^0(E)$ to the degeneracy locus of the evaluation morphism $W \otimes \mathcal{O}_X \xrightarrow{ev} E$, \cite{V1}. We may naturally identify the cokernel of $ev$ with $\mathrm{K}_C \otimes A^{-1}$, where $C=d(W) \in |L|$ and $A \in W^1_{k+1}(C)$, where $W^1_{k+1}(C)$ is the Brill--Noether locus of line bundles of degree $k+1$ with two sections, \cite{lazarsfeld-BNP}, \cite{aprodu-farkas}. If $C$ is sufficiently general, then $d^{-1}(C) \simeq W^1_{k+1}(C)$ is reduced, \cite{lazarsfeld-BNP}. Under the identification $d^{-1}(C) \simeq W^1_{k+1}(C)$, a line bundle $A \in W^1_{k+1}(C)$ is mapped to $\mathrm{H}^0(A)^{\vee}\seq \mathrm{H}^0(E)$. \smallskip
 
 
 By \cite[Proof of Prop.\ 7]{V1}, the spaces $\rm{Sym}^{k-2}\mathrm{H}^0(A)^{\vee}$, $A \in W^1_{k+1}(C)$ generate $\rm{Sym}^{k-2}\mathrm{H}^0(E)$. Each such $\mathrm{H}^0(A)^{\vee}$ corresponds to a line $T_A \seq \PP(\mathrm{H}^0(E))$ and set $\displaystyle{T:=\bigcup_{A \in W^1_{k+1}(C)} T_A}$ to be the union of these lines. The image of $T$ under $\widetilde{\psi}: \PP(\mathrm{H}^0(E)) \to \PP(\mathrm{K}_{k-1,1}(X,L))$ is non-degenerate so $\mathrm{K}_{k-1,1}(X,L)$ is generated by the subspaces $\mathrm{K}_{k-1,1}(X,L, \mathrm{H}^0(L \otimes I_{Z(s)})$, where $s \in \mathrm{H}^0(E)$ is a section corresponding to a $g^1_{k+1}$ on the fixed curve $C$. After restricting to $C$, such spaces are identified with subspaces of the form $$\mathrm{K}_{k-1,1}(C,\omega_C, \mathrm{H}^0(\omega_C \otimes A^{\vee})) \seq \mathrm{K}_{k-1,1}(C,\omega_C),\; \; A \in W^1_{k+1}(C)$$ under the isomorphism $\mathrm{K}_{k-1,1}(X,L) \simeq \mathrm{K}_{k-1,1}(C, \omega_C)$.
\end{proof}
